\documentclass[10pt,a4paper]{article}

% --- Core Packages ---
\usepackage[utf8]{inputenc}
\usepackage[T1]{fontenc}
\usepackage{lmodern}
\usepackage[margin=2.0cm, top=2.5cm, bottom=2.3cm, headheight=20pt]{geometry}
\usepackage{amsmath,amssymb}
\usepackage{booktabs}
\usepackage{tabularx}
\usepackage{xcolor}
\usepackage{colortbl}
\usepackage{enumitem}
\usepackage{titlesec}
\usepackage{fancyhdr}
\usepackage{tcolorbox}
\usepackage{microtype}
\usepackage[hidelinks, pdftex, colorlinks=true, linkcolor=navyblue, urlcolor=ceremateal, citecolor=navyblue]{hyperref}

% --- Custom Colors ---
\definecolor{ceremateal}{RGB}{0, 115, 150}
\definecolor{navyblue}{RGB}{0, 48, 87}
\definecolor{lightbg}{RGB}{246, 249, 252}
\definecolor{borderblue}{RGB}{180, 205, 225}
\definecolor{headergray}{RGB}{235, 242, 248}
\definecolor{mandatoryred}{RGB}{170, 30, 30}
\definecolor{optionalgray}{RGB}{100, 105, 110}

% --- Typography & Section Styling ---
\emergencystretch=1.5em

\titleformat{\section}
  {\Large\bfseries\color{navyblue}}{\thesection}{1em}{}[\color{ceremateal}\titlerule]
\titleformat{\subsection}
  {\large\bfseries\color{ceremateal}}{\thesubsection}{1em}{}

\titlespacing*{\section}{0pt}{1.2em}{0.6em}
\titlespacing*{\subsection}{0pt}{1.0em}{0.4em}

% --- Header and Footer ---
\pagestyle{fancy}
\fancyhf{}
\fancyhead[L]{\small\color{navyblue}\textbf{Cartino2D} – Architecture \& Software Dependencies}
\fancyhead[R]{\small\color{optionalgray}Cerema Hydraulics \& Hydrodynamics}
\fancyfoot[C]{\small\color{optionalgray}\thepage}
\renewcommand{\headrulewidth}{0.6pt}
\renewcommand{\headrule}{\hbox to\headwidth{\color{ceremateal}\leaders\hrule height \headrulewidth\hfill}}
\renewcommand{\footrulewidth}{0pt}

% --- Custom Box for Introduction ---
\tcbset{
  orchestratorbox/.style={
    colback=lightbg,
    colframe=borderblue,
    coltitle=navyblue,
    fonttitle=\bfseries,
    sharp corners,
    rounded corners=southeast,
    arc=4mm,
    boxrule=0.8pt,
    left=12pt,
    right=12pt,
    top=9pt,
    bottom=9pt
  }
}

\begin{document}

% --- Title Block ---
\begin{center}
  {\Huge\bfseries\color{navyblue} Cartino2D}\\[0.3em]
  {\Large\bfseries\color{ceremateal} Architecture \& Software Dependencies}\\[0.5em]
  {\small\color{optionalgray} Automated 2D Shallow-Water Hydrodynamic Modeling Framework – Cerema}\\[0.4em]
  \rule{\textwidth}{1.2pt}
\end{center}

\vspace{0.4em}

\begin{tcolorbox}[orchestratorbox]
\textbf{Cartino2D} is an automated framework developed by \textbf{Cerema} to design, execute, and post-process 2D shallow-water hydrodynamic simulations (fluvial flooding, direct pluvial runoff, and marine submersion).

In this architecture, \textbf{R acts as the central driver and orchestrator (the ``pilot'')}. It does not solve the heavy computational fluid dynamics directly; instead, it coordinates a specialized ecosystem of open-source numerical solvers, mesh generators, GIS engines, and high-performance computing (HPC) utilities.
\end{tcolorbox}

\vspace{0.6em}

\section*{1. Summary of Cartino2D Dependencies}
\addcontentsline{toc}{section}{1. Summary of Cartino2D Dependencies}

\noindent
\begin{tabularx}{\textwidth}{>{\bfseries\raggedright\arraybackslash}p{2.6cm} X >{\raggedright\arraybackslash}p{3.2cm} >{\raggedright\arraybackslash}p{2.5cm} >{\centering\arraybackslash}p{2.1cm}}
\toprule
\rowcolor{headergray}
\color{navyblue}Component & \color{navyblue}Core Role & \color{navyblue}Language / Runtime & \color{navyblue}Cartino2D Step & \color{navyblue}Status \\
\midrule
R \& RStudio & Master Driver \& Orchestrator & R (4.4+) & All Steps (00--09, PreC2D) & \textbf{\color{mandatoryred}Mandatory} \\
\addlinespace[3pt]
TELEMAC-2D & 2D Hydrodynamic numerical solver & Fortran / Python / MPI & Step 05 (Local \& Remote) & \textbf{\color{mandatoryred}Mandatory} \\
\addlinespace[3pt]
GMSH & 2D Unstructured finite-element mesher & C++ (64-bit binary) & Step 02 (Mesh generation) & \textbf{\color{mandatoryred}Mandatory} \\
\addlinespace[3pt]
PPutils & Selafin \texttt{.slf} creation \& result raster extraction & Python & Steps 01, 02 \& Step 07 & \textbf{\color{mandatoryred}Mandatory} \\
\addlinespace[3pt]
Python & Runtime engine for PPutils scripts & Python 3 (\texttt{numpy}, \texttt{scipy}) & Steps 01, 02, 07 & \textbf{\color{mandatoryred}Mandatory} \\
\addlinespace[3pt]
GRASS GIS & High-performance raster engine \& culvert repair & GIS CLI (\texttt{grass84.bat}) & Step 01 \& Step 08 & \textbf{\color{mandatoryred}Mandatory} \\
\addlinespace[3pt]
QGIS \& OSGeo4W & Headless GIS processing, GDAL utilities, styles \& \texttt{.gpkg} & GIS CLI (\texttt{qgis\_process}) & PreC2D, 00, 01, 02, 07--09 & \textbf{\color{mandatoryred}Mandatory} \\
\addlinespace[3pt]
TauDEM & Hydrological terrain analysis (DEM conditioning) & C++ binaries (\texttt{.exe}) / MS-MPI & PreC2D & \textbf{\color{mandatoryred}Mandatory} \\
\addlinespace[3pt]
PSCP (PuTTY) & SSH/SCP encrypted bridge to HPC clusters & Windows CLI executable & Step 05 (Remote HPC runs) & \color{optionalgray}Optional (HPC) \\
\addlinespace[3pt]
Fortran Routines & Custom user laws (rain, hydrograph interpolation) & Fortran 90 & Steps 04 / 05 & Built-in C2D \\
\addlinespace[3pt]
Rtools & C/C++/Fortran Windows compiler toolchain & Windows build tools & R geospatial compilation & \textbf{\color{mandatoryred}Mandatory} \\
\bottomrule
\end{tabularx}

\vspace{1.2em}

\section*{2. Detailed Component Descriptions}
\addcontentsline{toc}{section}{2. Detailed Component Descriptions}

\subsection*{2.1 R \& RStudio — The Master Orchestrator (The Driver)}
\textbf{R} is the backbone and brain of Cartino2D. Instead of executing heavy computational fluid dynamics directly in R, it acts as an automation framework:
\begin{itemize}[leftmargin=1.5em, itemsep=2pt, topsep=2pt]
  \item Reads user-defined spatial extent, grid resolutions, and boundary condition files.
  \item Automates domain sectorization and generates all intermediate configuration files (\texttt{.cas}, \texttt{.cli}, \texttt{.geo}, \texttt{.bat}) required by downstream software.
  \item Spawns background system processes and coordinates multicore execution using R packages (\texttt{parallel}, \texttt{future}, \texttt{sf}, \texttt{terra}).
  \item Manages error logging, workspace initialization, and folder structures.
\end{itemize}

\subsection*{2.2 TELEMAC-2D \& MS-MPI — The 2D Hydrodynamic Numerical Solver}
Developed by the Open TELEMAC-MASCARET international consortium, \textbf{TELEMAC-2D} solves the two-dimensional shallow-water Saint-Venant equations on unstructured triangular meshes.
\begin{itemize}[leftmargin=1.5em, itemsep=2pt, topsep=2pt]
  \item Solves water depths ($h$), depth-averaged horizontal velocity fields ($u, v$), bottom friction (Strickler laws), and dynamic infiltration (SCS-CN model).
  \item \textbf{Natively computes Control Section (SC) balances}: During the simulation, TELEMAC-2D internally integrates water fluxes across user-defined control sections at each time step, outputting instantaneous discharges and cumulative volumes directly into the simulation listings (\texttt{sections\_output.txt} / output listings). R directly parses these text logs in Step 06 to plot interactive hydrographs.
  \item Simulates river flooding, pluvial surface runoff, and coastal submersion.
  \item Leverages \textbf{Microsoft MPI (MS-MPI)} under Windows to split meshes across multiple CPU cores for high-speed parallel computation.
\end{itemize}

\subsection*{2.3 GMSH — 2D Unstructured Finite-Element Mesher}
\textbf{GMSH} is an open-source 3D/2D finite-element mesh generator with a built-in CAD engine.
\begin{itemize}[leftmargin=1.5em, itemsep=2pt, topsep=2pt]
  \item Cartino2D feeds boundary contours, elevation breaklines (embankments, roads, ditches, valley bottoms), and variable target mesh sizes into GMSH.
  \item GMSH generates a high-resolution, topography-adaptive Delaunay triangulation (\texttt{.msh}).
  \item R immediately translates this mesh into the standard TELEMAC binary \textbf{Selafin (\texttt{.slf})} format.
\end{itemize}

\subsection*{2.4 PPutils \& Python — Mesh Geometry \& Selafin Assembly / Extraction}
\textbf{PPutils} (\textit{Post-Processing utilities}) is an open-source Python toolchain adapted and enriched by Cerema for hydraulic engineering. It intervenes at multiple key steps of the Cartino2D modeling chain:
\begin{itemize}[leftmargin=1.5em, itemsep=2pt, topsep=2pt]
  \item \textbf{Pre-simulation (Steps 01 \& 02 - Selafin Creation)}: Intervenes after GMSH meshing. Using scripts such as \texttt{append\_adcirc\_cerema.py}, it fuses the triangular mesh grids (\texttt{.grd}), interpolated bathymetry/elevation, spatially distributed roughness (Strickler friction), and soil infiltration capacity (Curve Number) into the official binary \textbf{Selafin geometry file (\texttt{.slf})} and boundary conditions file (\texttt{.cli}) required by TELEMAC-2D.
  \item \textbf{Post-simulation (Step 07 - Mesh-to-Raster Conversion)}: Decodes massive binary TELEMAC-2D mesh simulation results (\texttt{.slf} / \texttt{.res}) and converts peak hydrodynamic envelopes into standard ESRI ASCII raster grids (\texttt{.asc} via \texttt{slf2asc.py}) for water depths (\texttt{HWH}), velocities (\texttt{HV}), unit discharges (\texttt{HSF}), and bathymetry (\texttt{Bottom\_m}).
\end{itemize}
\textit{Note: PPutils does not calculate control section discharges; those are computed natively by TELEMAC-2D.}

\subsection*{2.5 GRASS GIS — High-Performance Raster Engine \& Culvert Repair}
\textbf{GRASS GIS} is invoked in headless command-line mode (\texttt{grass84.bat}). It handles heavy raster operations without the memory overhead of desktop GIS applications:
\begin{itemize}[leftmargin=1.5em, itemsep=2pt, topsep=2pt]
  \item \textbf{Spatial masking \& cleaning}: Sets residual water depths to null (\texttt{r.null}) below physical thresholds (e.g., $< 0.005\,\text{m}$).
  \item \textbf{Hydraulic culvert anti-hole restoration (Anti-trous)}: During meshing, hydraulic structures (bridges, culverts) are often carved into the 1D/2D mesh. In Step 08, GRASS uses map algebra (\texttt{r.mapcalc}) to restore the natural ground elevation (\texttt{MNT\_BRUT\_NAT}) around structures, preventing artificial sinkholes from biasing final inundation maps.
\end{itemize}

\subsection*{2.6 QGIS \& OSGeo4W (\texttt{qgis\_process}) — Headless Vector/Raster GIS \& Cartography}
Cartino2D harnesses \textbf{QGIS} (in headless batch mode via \texttt{qgis\_process}) and the \textbf{OSGeo4W} GDAL toolchain across multiple phases of the pipeline, far beyond just post-processing:
\begin{itemize}[leftmargin=1.5em, itemsep=2pt, topsep=2pt]
  \item \textbf{PreC2D \& Step 00 (Data Preparation \& Sectorization)}: Executes geometric vector operations (polygon buffering, differences, layer merges, deleting holes via \texttt{native:deleteholes}, and selecting main connected components) as well as raster clipping and format translations.
  \item \textbf{Step 01 \& Step 02 (Topography \& Mesh Verification)}: Builds virtual elevation mosaics (\texttt{gdalbuildvrt}), generates overviews (\texttt{gdaladdo}), and inspects 2D mesh topology via \texttt{native:exportmeshfaces} before simulation launch.
  \item \textbf{Steps 07, 08 \& 09 (Post-Processing \& Delivery)}: Performs mesh rasterization (\texttt{native:meshrasterize}), reprojections (\texttt{gdal:warpreproject}), crops rasters to validated zones, automatically injects official Cerema symbology templates (\texttt{.qml}), and packages multi-scenario outputs into compressed \textbf{OGC GeoPackage (\texttt{.gpkg})} databases with built-in pyramids.
\end{itemize}

\subsection*{2.7 TauDEM — Hydrological Terrain Conditioning (PreC2D)}
\textbf{TauDEM} (\textit{Terrain Analysis Using Digital Elevation Models}) is a \textbf{mandatory} suite of compiled standalone C++ command-line executables (\texttt{.exe}) developed by Utah State University. It runs in parallel via \textbf{Microsoft MPI (\texttt{mpiexec})} and does \textbf{not} rely on Python:
\begin{itemize}[leftmargin=1.5em, itemsep=2pt, topsep=2pt]
  \item Removes artificial terrain depressions and sinks from raw LiDAR/DEM data (\texttt{PitRemove.exe}).
  \item Computes flow directions (\texttt{D8Flowdir.exe} and \texttt{DinfFlowDir.exe}) and contributing drainage areas (\texttt{AreaD8.exe}).
  \item Delineates sub-catchments and stream networks to guide the placement of rainfall and liquid boundary sources.
\end{itemize}

\subsection*{2.8 PSCP (PuTTY) — Encrypted Supercomputing Bridge (HPC)}
For massive regional or national simulations that exceed the capacity of local desktop workstations, \textbf{\texttt{pscp.exe}} provides an automated, secure SCP transfer pipeline:
\begin{itemize}[leftmargin=1.5em, itemsep=2pt, topsep=2pt]
  \item Uploads geometry meshes (\texttt{.slf}), steering files (\texttt{.cas}), and boundary condition files (\texttt{.cli}) to remote high-performance supercomputers (e.g., IFREMER Datarmor, CEA TGCC Irene, DTecRem Margny).
  \item Synchronizes and downloads result files (\texttt{.res}, \texttt{.slf}) back to the local workstation once Slurm batch jobs complete.
\end{itemize}

\subsection*{2.9 Dynamic Fortran Routines (\texttt{.f}) — Advanced Hydraulic Laws}
Located in Cartino2D subfolders, custom Fortran 90 routines are compiled on the fly by TELEMAC at run launch:
\begin{itemize}[leftmargin=1.5em, itemsep=2pt, topsep=2pt]
  \item \texttt{discharge\_interpol.f}: Performs dynamic temporal linear interpolation of discharge hydrographs injected at inflow boundaries and liquid sources.
  \item Spatially and temporally distributes radar rainfall matrices and soil infiltration laws across mesh nodes.
\end{itemize}

\subsection*{2.10 Rtools — Windows Native Build Toolchain}
\textbf{Rtools} provides the GCC, Gfortran, and C++ compilers required on Windows:
\begin{itemize}[leftmargin=1.5em, itemsep=2pt, topsep=2pt]
  \item Allows R to compile native C/C++ dependencies behind core geospatial packages (\texttt{sf}, \texttt{terra}, \texttt{units}, \texttt{gdalraster}).
  \item Ensures binary compatibility and high-performance execution of geospatial algebra in R.
\end{itemize}

\end{document}
